\Article{12}
Als viertes Beispiel wollen wir die Abbildung der Oberfläche eines Rotationsellipsoids in der Ebene betrachten. Seien \(a\) und \(b\) die beiden Haupthalbachsen des Ellipsoids, sodass sich die Koordinaten durch
\[
\begin{aligned}
x&=a\cos t\sin u\\
y&=a\sin t\sin u\\
z&=b\cos u
\end{aligned}
\]
\SourceLeaf{211}%
\noindent angeben lassen. Daraus ergibt sich
\[
\omega=aa\sin u^2\,\mathrm{d}t^2+(aa\cos u^2+bb\sin u^2)\mathrm{d}u^2
\]
\noindent und die Differentialformel \(\omega=0\) liefert, wenn wir der Kürze halber \(\BareRoot{(1-\frac{bb}{aa})}=\varepsilon\) setzen, sofern für die Halbachse entlang der Rotationsachse \(b<a\) gilt,
\[
0=\mathrm{d}t\mp i\,\mathrm{d}u\mathord{.}\BareRoot{(\mathrm{cotang}\,u^2+1-\varepsilon\varepsilon)}
\]

Setzt man hier
\[
\BareRoot{(1-\varepsilon\varepsilon)}\mathord{.}\mathrm{tang}\,u=\mathrm{tang}\,w
\]
\noindent wobei bei der Anwendung auf das Erdsphäroid \(90^0-w\) die geografische Breite und \(t\) die Länge bezeichnet, so nimmt diese Gleichung die Form
\[
0=\mathrm{d}t\mp i\,\mathrm{d}w\mathord{.}\frac{1-\varepsilon\varepsilon}{(1-\varepsilon\varepsilon\cos w^2)\sin w}
\]
\noindent an. Durch Integration erhält man
\[
\text{Const.}=t\pm i\log\mathord{.}\Bigl\{\mathrm{cotang}\tfrac{1}{2}w\mathord{.}\Bigl(\frac{1-\varepsilon\cos w}{1+\varepsilon\cos w}\Bigr)^{\frac{1}{2}\varepsilon}\Bigr\}
\]
\noindent Ist \(f\) eine beliebige Funktion, so nimmt man daher für \(X\) den Realteil und für \(iY\) den Imaginärteil von
\[
f\Bigl(t+i\log\Bigl\{\mathrm{cotang}\tfrac{1}{2}w\mathord{.}\Bigl(\frac{1-\varepsilon\cos w}{1+\varepsilon\cos w}\Bigr)^{\frac{1}{2}\varepsilon}\Bigr\}\Bigr)
\]
\noindent --- Wählt man für \(f\) eine lineare Funktion, also \(f\upsilon=k\upsilon\), so erhält man
\[
X=kt,\quad Y=k\log\mathrm{cotang}\tfrac{1}{2}w-\tfrac{1}{2}k\varepsilon\log\frac{1+\varepsilon\cos w}{1-\varepsilon\cos w}
\]
\noindent und damit eine zur \textsc{Mercator}-Projektion analoge Projektion.

Wählt man dagegen für \(f\) eine imaginäre Exponentialfunktion \(f\upsilon=ke^{i\lambda\upsilon}\), so erhält man
\[
\begin{aligned}
X&=k\mathord{.}\mathrm{tang}\tfrac{1}{2}w^{\lambda}\mathord{.}\Bigl(\frac{1+\varepsilon\cos w}{1-\varepsilon\cos w}\Bigr)^{\frac{1}{2}\varepsilon\lambda}\mathord{.}\cos\lambda t,\\
Y&=k\,\mathrm{tang}\tfrac{1}{2}w^{\lambda}\mathord{.}\Bigl(\frac{1+\varepsilon\cos w}{1-\varepsilon\cos w}\Bigr)^{\frac{1}{2}\varepsilon\lambda}\mathord{.}\sin\lambda t
\end{aligned}
\]
\noindent Setzt man \(\lambda=1\), so ergibt sich eine zur stereografischen Polarprojektion analoge Projektion. Im Allgemeinen ergibt sich eine sehr zweckmäßige Projektion zur Abbildung eines Teils der Erdoberfläche, wenn ihre Abplattung berücksichtigt werden soll.

Was über den anderen Fall mit \(b>a\) zu sagen ist, lässt sich zwar leicht unmittelbar aus dem Vorhergehenden ableiten. Behält man dieselben
\SourceLeaf{212}
Bezeichnungen bei, so wird \(\varepsilon\) imaginär, aber \((\frac{1+\varepsilon\cos w}{1-\varepsilon\cos w})^{\frac{1}{2}\varepsilon}\) wieder reell. Der Vollständigkeit wegen wollen wir die Formeln für diesen Fall jedoch eigens angeben und gleich zu Beginn \(\BareRoot{(\frac{bb}{aa}-1)}=\eta\) setzen. Dann ist \(w\) aus der Gleichung
\[
\BareRoot{(1+\eta\eta)}\mathord{.}\mathrm{tang}\,u=\mathrm{tang}\,w
\]
\noindent zu bestimmen, und die Differentialgleichung
\[
0=\mathrm{d}t\mp i\,\mathrm{d}w\mathord{.}\frac{1+\eta\eta}{(1+\eta\eta\cos w^2)\sin w}
\]
\noindent liefert das Integral
\[
\text{Const.}=t\pm i(\log\mathrm{cotang}\tfrac{1}{2}w+\eta\,\mathrm{Arc}\,\mathrm{tang}.\,\eta\cos w)
\]
\noindent sodass man \(X\) mit dem Realteil und \(iY\) mit dem Imaginärteil von
\[
f(t+i(\log\mathrm{cotang}\tfrac{1}{2}w+\eta\,\mathrm{Arc}\,\mathrm{tang}.\,\eta\cos w))
\]
\noindent gleichsetzen muss. Die Gegenstücke zu den beiden obigen speziellen Anwendungen ergeben sich unmittelbar daraus. Bei der ersten ist
\[
X=kt,\qquad Y=k\log\mathrm{cotang}\tfrac{1}{2}w+\eta k\,\mathrm{Arc}\,\mathrm{tang}.\,\eta\cos w
\]
\noindent und bei der anderen
\[
\begin{aligned}
X&=k\,\mathrm{tang}\tfrac{1}{2}w^{\lambda}\mathord{.}e^{-\eta\lambda\,\mathrm{Arc}\,\mathrm{tang}.\,\eta\cos w}\mathord{.}\cos\lambda t\\
Y&=k\,\mathrm{tang}\tfrac{1}{2}w^{\lambda}\mathord{.}e^{-\eta\lambda\,\mathrm{Arc}\,\mathrm{tang}.\,\eta\cos w}\mathord{.}\sin\lambda t
\end{aligned}
\]
\noindent zu setzen.

